\section{Frameworks, standards, and industry practice}\label{sec:standards}

\subsection{OWASP Top 10 for LLM Applications: 2023 $\to$ 2026}\label{sec:owasp}

\begin{table}[htbp]
\centering
\small
\caption{The OWASP list across its three editions (abbreviated names). The 2026
edition is the agentic turn: acting systems, cascading effects, cross-system
trust.}\label{tab:owasp-evolution}
\begin{tabular}{L{4.5cm} L{4.3cm} L{4.5cm}}
\toprule
2023 (LLM01--10) & 2025 & 2026 (agentic edition) \\
\midrule
Prompt injection & Prompt injection (kept \#1) & Agentic directive injection
family; hidden context exposure replaces system-prompt leakage \\
Insecure output handling & Sensitive information disclosure & Excessive agency
rises to \#3---tool effects, cascading actions \\
Training data poisoning & Supply chain & Tool/MCP ecosystem classes first-class \\
Plugin misuse & Data and model poisoning & Cross-system trust and second-order
output chains \\
Excessive agency & Improper output handling & Persistent memory classes
recognized \\
Unbounded consumption / leaky APIs & Unbounded consumption & Availability and
cost-amplification classes consolidated; persistent memory recognized \\
\bottomrule
\end{tabular}
\end{table}

Three readings of the reshuffle (expanded in Section~\ref{sec:owasp2026}): the
unit of risk became the acting system; the delivery channel became retrieval
itself; and persistence (memory, cross-run state) entered the threat model as a
first-class citizen. Our taxonomy (Table~\ref{tab:taxonomy}) is deliberately a
superset: evaluation-validity classes (W7--W9) are not OWASP categories at all,
because they are weaknesses in \emph{us}, the measurers.

\subsection{MITRE ATLAS, NIST AI RMF, EU AI Act, ISO 42001}\label{sec:standards-map}

MITRE ATLAS (data release v2026.07) carries 101 techniques and 77
sub-techniques grounded in 68 case studies, with the AML.T0051 family
tagging prompt injection against agentic systems \cite{atlas}. For engineering use, ATLAS's value is the mapping discipline:
every technique we model (Section~\ref{sec:taxonomy}) should carry an ATLAS
anchor so that telemetry, detections, and reports speak a shared language.

NIST's AI RMF plus its Generative AI Profile supplies the governance skeleton:
\emph{measure} and \emph{manage} functions map almost one-to-one onto our cell
cards and remediation queue \cite{nist-ai-rmf,nist-genai}. The EU AI Act adds
dates: entered into force August 2024; GPAI obligations from August 2025; the
high-risk and legacy-system milestones run through August 2026--2027
\cite{eu-ai-act}. ISO/IEC 42001 gives the management-system wrapper
certification audits will ask for \cite{iso42001}. None of these three tells an
engineer \emph{what to build}; they tell an operator \emph{what to prove} ---
which is why this paper's statistics exist: to make the proving honest.

\subsection{Industry red-team tooling and the missing layer}\label{sec:tooling}

The mature open tooling covers static probe execution: PyRIT (Microsoft's
automation pipeline), garak (probe suites), promptfoo (eval harnesses), plus
bug-bounty programs that have formalized AI scopes. Two gaps persist. First,
\emph{adaptive measurement}: most suites are one-shot; the escalation ladder and
guard-wear reporting of Section~\ref{sec:gaps} are rare. Second,
\emph{honest statistics}: suites report rates, rarely intervals, almost never a
retry-inflation null. The statistical layer of Section~\ref{sec:stats} is our
fill for that gap, and it is deliberately a library, not a product: any cell
card in any framework should be able to import it.

\subsection{The commercial layer: who pays, and what it teaches}\label{sec:money}

Between September 2025 and March 2026 the tooling landscape consolidated
faster than the literature could record it. OpenAI agreed to acquire
promptfoo (March~9, 2026), folding the most widely deployed open-source
red-team harness---350{,}000+ developers reached, 130{,}000 monthly
active, used by more than a quarter of the Fortune~500---into its
Frontier platform while pledging to keep the core open source
\cite{promptfoo-openai}. Check Point acquired Lakera (announced
September~2025; terms undisclosed, press estimates $\approx$\$300M),
pairing Lakera Guard's runtime enforcement with Lakera Red's automated
testing and the Gandalf research funnel's 80M+ recorded attacks
\cite{checkpoint-lakera,lakera-ctech}. SentinelOne acquired Prompt
Security for GenAI runtime defense (August~2025)
\cite{sentinelone-promptsec}. Microsoft left PyRIT open but sells its
capability as the Foundry AI Red Teaming Agent \cite{foundry-rt}; garak
remains free under NVIDIA. HackerOne monetizes the opposite direction:
human creativity at bounty scale---a 540\% year-over-year surge in
validated prompt-injection reports in 2025, and a human-led,
agent-driven AI Red Teaming service \cite{hackerone-hpsr,hackerone-ai-rt}.

Five revenue models now define the market:
\begin{itemize}
  \item \textbf{Runtime protection} (Lakera Guard, Prompt Security):
  priced per protected application and call volume; the classic
  security-budget line.
  \item \textbf{Continuous red-team SaaS} (Lakera Red, the Foundry
  agent, HackerOne AI Red Teaming): subscriptions for scheduled
  adversarial campaigns with remediation workflows.
  \item \textbf{Developer-platform probe volume} (promptfoo): free
  community tier up to 10k probes per month, enterprise pricing for
  scale, SSO, and compliance dashboards \cite{promptfoo-pricing}.
  \item \textbf{Human bounty markets} (HackerOne): platform fees on
  researcher payouts, now the largest single channel of AI finding
  volume.
  \item \textbf{Compliance-driven audit} (EU AI Act Art.~55 evidence,
  ISO/IEC 42001): sold today as manual services; the least automated
  of the five.
\end{itemize}

For this project the lesson is positioning, not envy. First, platform
vendors buying the standard harness sharpens the auditor-independence
problem: a model vendor grading its own exam, exactly the structural
conflict regulated buyers already refuse in finance. An independent,
statistics-first verdict layer is the open wedge. Second, every one of
these products sells finding \emph{counts}; none sells calibrated
\emph{uncertainty}. The statistical layer of Section~\ref{sec:stats}
maps directly onto the compliance-audit model---the least automated and,
given the 2026--2027 EU milestones, the fastest-moving of the five.

\subsection{Incident record}\label{sec:incidents}

The incident record is the taxonomy's reality check: every entry below maps
to at least one W-class, and in no case did vendor-side monitoring surface
the failure before researchers, customers, or viral exposure did.

\begin{table}[htbp]
\centering
\small
\caption{Verified incidents and advisories that anchor the taxonomy's realism.}
\label{tab:incidents}
\begin{tabular}{L{3.2cm} L{5.0cm} L{5.2cm}}
\toprule
Case & Class & Outcome / lesson \\
\midrule
EchoLeak, M365 Copilot (CVE-2025-32711) \cite{echolake-cve,echolake-paper} &
W6 + W1 & Zero-click exfil via email-seeded injection; reference-style Markdown
bypassed redaction; image auto-fetch carried data out. \\
MCP Inspector RCE (CVE-2025-49596) \cite{cve-2025-49596,oligo-mcp} & W4 &
Developer tooling itself was the entry point; CVSS 9.4; fixed in 0.14.1. \\
Cursor / MCP-STDIO family (CVE-2025-54136 et al.) \cite{ox-mcp-family} & W4 &
Same STDIO injection pattern across vendors; supply-chain class is systemic. \\
Air Canada chatbot (2024 tribunal) & W12 & Airline bound by its chatbot's
promise; output chains carry legal weight. \\
DPD chatbot swearing (2024) & W12 & Output-control failure; customer-facing
reputational damage. \\
Chevrolet of Watsonville ``\$1 Tahoe'' (2025) & W12 + W11 & Agent-shaped
commercial surface accepted an absurd commitment; the offer went viral. \\
\bottomrule
\end{tabular}
\end{table}

\subsection{Books}\label{sec:books}

Two books anchor the craft literature. \emph{Red Teaming} (Siva Kumar and
Anderson, eds.) collects the enterprise-era practices: threat modeling for ML,
abuse-case elicitation, and the operational choreography of AI red teams
\cite{sivakumar-anderson}. Sotiropoulos's \emph{AI Penetration Testing} is the
practitioner's field manual: recon through exploitation on LLM stacks with
working toolchains \cite{sotiropoulos}. Both predate the 2026 agentic reshuffle,
which is precisely why this paper exists: the gap they leave---agents acting,
persisting, and propagating---is what Sections~\ref{sec:agentic}
and~\ref{sec:gaps} fill.
